%% Elsevier elsarticle template — numbered references
%% Journal of Hazardous Materials  /  Science of the Total Environment
%%
%% Uncomment the preferred layout option below.
%% For review (double-spaced): \documentclass[authoryear,preprint,review,12pt]{elsarticle}

%\documentclass[final,5p,times,twocolumn]{elsarticle}
%% \documentclass[final,5p,times]{elsarticle}          % single-column 5p
%% \documentclass[final,3p,times,twocolumn]{elsarticle} % two-column 3p
 \documentclass[preprint,12pt]{elsarticle}            % preprint / review

%% -------------------------------------------------------------------------
%% Packages
%% -------------------------------------------------------------------------
\usepackage{amssymb}
\usepackage{amsmath}
\usepackage{graphicx}
\graphicspath{{figures/}}
\usepackage{multirow}
\usepackage{booktabs}
\usepackage{tabularx}
\usepackage{xcolor}
\usepackage{float}
\usepackage{adjustbox}
\usepackage{lineno}          % line numbers for review — activate with \linenumbers

%% New centred tabularx column type
\newcolumntype{Y}{>{\centering\arraybackslash}X}

%% -------------------------------------------------------------------------
%% Journal name (change to match target journal)
%% -------------------------------------------------------------------------
\journal{Journal of Hazardous Materials}

%% =========================================================================
\begin{document}
%% =========================================================================

\begin{frontmatter}

%% -------------------------------------------------------------------------
%% Title
%% -------------------------------------------------------------------------
\title{AI-Driven Soil Pollution Assessment Through Computational and
    Spectroscopic Methodologies}

%% -------------------------------------------------------------------------
%% Authors and affiliations
%% -------------------------------------------------------------------------
\author[aff1,aff2,aff3]{Polina Lemenkova\corref{cor1}}
\ead{polina.lemenkova@tech-orleans.fr}
%% ORCID: 0000-0002-5759-1089

\cortext[cor1]{Corresponding author.}

\affiliation[aff1]{organization={Bureau de Recherches
    G\'eologiques et Mini\`eres (BRGM)},
    addressline={3 Avenue Claude-Guillemin},
    city={Orl\'eans Cedex~2},
    postcode={45060},
    country={France}}

\affiliation[aff2]{organization={Institut d'Etudes Avanc\'ees du
    Val de Loire, LE STUDIUM},
    addressline={1 Rue Dupanloup},
    city={Orl\'eans},
    postcode={45000},
    country={France}}

\affiliation[aff3]{organization={La Technopole d'Orl\'eans},
    addressline={1 Avenue du Champ de Mars},
    city={Orl\'eans},
    postcode={45100},
    country={France}}

%% -------------------------------------------------------------------------
%% Abstract
%% -------------------------------------------------------------------------
\begin{abstract}
Soil pollution has emerged as a critical global environmental challenge driven
by rapid industrialization, intensive agriculture, urban expansion, mining
activities, and the excessive use of fertilizers and pesticides. This review
presents an integrated framework for soil pollution assessment with emphasis
on major pollutant groups, including heavy metals, pesticides, nutrients,
microplastics, and per- and polyfluoroalkyl substances (PFAS). Drawing on
evidence from more than 100 peer-reviewed studies, the paper synthesizes
current advances in computational, spectroscopic, and artificial intelligence
(AI)-driven methodologies for soil monitoring and environmental risk
evaluation. Traditional analytical techniques based on extraction and
laboratory measurements are compared with advanced spectroscopic approaches,
including Fourier transform infrared spectroscopy (FTIR), X-ray photoelectron
spectroscopy (XPS), hyperspectral imaging, and synchrotron radiation
technologies, which enable precise characterization of pollutant speciation
and spatial distribution within soil aggregates. Recent developments in AI,
machine learning, computer vision, and remote sensing have significantly
enhanced the detection, classification, and prediction of soil contaminants
through process-based, empirical, composite, and statistical modelling
frameworks. Techniques such as Principal Component Analysis (PCA), Monte
Carlo simulations, geochemical modelling, and time-series forecasting provide
improved capabilities for identifying pollution hotspots, predicting
contaminant transport, and evaluating ecological risks to soil ecosystem
services (SES). Despite these advances, major challenges remain related to
heterogeneous soil properties, limited detection sensitivity, data
uncertainty, and model complexity. Future research should prioritize the
integration of AI-driven analytics, real-time sensor technologies, remote
sensing platforms, and predictive environmental modelling to support
sustainable soil management, ecosystem protection, and evidence-based
environmental policy development.
\end{abstract}

%% -------------------------------------------------------------------------
%% Graphical abstract (optional — insert figure path or comment out)
%% -------------------------------------------------------------------------
\begin{graphicalabstract}
%% \includegraphics[width=\textwidth]{graphical_abstract}
\end{graphicalabstract}

%% -------------------------------------------------------------------------
%% Research highlights
%% -------------------------------------------------------------------------
\begin{highlights}
\item Integrative review of AI and ML methodologies for soil pollution
    assessment across five major pollutant groups.
\item Synchrotron, FTIR, and XPS spectroscopic techniques are benchmarked
    against AI-driven analytical frameworks.
\item A six-step multi-stage workflow for SES modelling and ecological risk
    quantification is proposed.
\item Physics-Informed Neural Networks (PINNs) and ensemble learning offer
    superior contaminant transport predictions.
\item Future priorities include real-time sensor integration, explainable AI,
    and standardized monitoring frameworks.
\end{highlights}

%% -------------------------------------------------------------------------
%% Keywords
%% -------------------------------------------------------------------------
\begin{keyword}
soil pollution \sep artificial intelligence \sep machine learning \sep
spectroscopic methods \sep heavy metals \sep PFAS \sep soil ecosystem services
\end{keyword}

\end{frontmatter}

%% \linenumbers   % uncomment for review submission

%% =========================================================================
\section{Introduction}\label{sec1}
%% =========================================================================

%% -------------------------------------------------------------------------
\subsection{Background}\label{sec1.1}
%% -------------------------------------------------------------------------

Soil ecosystem services (SES) play a fundamental role in sustaining
environmental quality and supporting human well-being through the regulation
of soil functions and vegetation dynamics. Effective management of SES
requires continuous assessment of climatic and hydrological conditions
\cite{Lal2020,Robinson2014}, optimization of agricultural practices
\cite{Hepperly2009,Rahman2014}, regulation of nutrient cycling and organic
matter decomposition \cite{Atoloye2024,Ma2024,Vishwanath2025}, and early
identification of soil contamination
\cite{Pan2019,Fan2025,Lemenkova2025b,Proshad2025}. Soil degradation is
increasingly associated with the accumulation of diverse pollutants,
including pesticides (e.g., insecticides, herbicides, fungicides, and
bactericides), mineral oils, microplastics, per- and polyfluoroalkyl
substances (PFAS), excess nutrients, and heavy metals
\cite{Shen2024,Guan2024,Baghaie2020,Luo2017}. The growing spatial extent of
contaminated soils represents a major environmental and ecological concern,
emphasizing the urgent need for robust and integrated monitoring strategies
\cite{Liu2020,Zhang2021,Lemenkova2025d}. Recent studies have demonstrated
that soil pollution is closely linked to both natural hazards and
anthropogenic activities, particularly industrialization
\cite{Rosskopfova2012,Guo2025,Zhang2022} and intensive agriculture
\cite{Hepperly2009,Rahman2014,Pan2019}, which contribute substantially to the
accumulation and redistribution of contaminants in terrestrial ecosystems.

Advances in computational and analytical frameworks have revolutionized the
identification of soil constituents and their interactions with pollutants
\cite{Liu2020,Zhang2021}. Modern ecological monitoring utilizes spectroscopic
analysis, machine learning (ML), and remote sensing to quantify pollutant
transport, transformation, and persistence across scales
\cite{Fan2025,Lemenkova2025d}. Specifically, integrating deep learning
algorithms with hyperspectral imagery enables high-fidelity characterization
of contaminant dynamics within organic-mineral matrices \cite{Xu2026,Li2023}.
These predictive frameworks optimize risk assessment and environmental
management by capturing complex multi-temporal spatial heterogeneities
\cite{Proshad2025,Feng2024}.

Modern computational architectures have fundamentally transformed the
modelling of soil contaminant kinetics, facilitating high-fidelity simulations
of risk to ecosystem services
\cite{Liu2020,Zhang2021,Proshad2025,Lemenkova2025d}. While traditional
identification frameworks rely on process-based
\cite{Giannakis2017,Karlen2019}, empirical \cite{Benavidez2018,Lal2020},
composite \cite{Sadhukhan2022}, and statistical models
\cite{Lu2020,Feng2024,Lemenkova2024b}, the field is rapidly pivoting toward
Artificial Intelligence (AI) and Machine Learning (ML) to handle complex
geochemical data. Integration of AI-driven computational methods
\cite{Proshad2025,Lemenkova2025d,Feng2024,Li2023} with hyperspectral imaging
\cite{Fan2025,Shi2021}, spectroscopic sensors \cite{Xu2026,Guo2025}, and
remote sensing \cite{Lemenkova2025b,Zhang2022} allows for the automated
extraction of non-linear contaminant patterns across heterogeneous soil
systems. These ML architectures---specifically deep learning (DL) and
ensemble classifiers---excel where standard models struggle with data
uncertainty \cite{Lu2020,Karlen2019,Li2023} and instrumental sensitivity
\cite{Rosskopfova2012}. Despite challenges regarding model transferability
\cite{Benavidez2018,Zhang2021} and soil complexity
\cite{Klaucol2017,Lemenkova2023b}, the priority has shifted toward scalable,
data-driven pipelines. By merging neural networks and geospatial analytics,
researchers can achieve real-time predictive monitoring, drastically improving
environmental risk assessment and the long-term protection of soil ecosystems.

%% -------------------------------------------------------------------------
\subsection{Research Problem}\label{sec1.2}
%% -------------------------------------------------------------------------

The inherent spatial heterogeneity of soil microstructure dictates the
stochastic mobility and kinetic transformation of contaminants, requiring
sophisticated modelling to map these multi-scale interactions
\cite{Vogel2018}. Stochastic modelling of soil systems accounts for the
inherent randomness and multiscale variability of porous media, where fluid
and solute transport are governed by probabilistic frameworks rather than
linear trajectories \cite{Vogel2018,Lemenkova2023}. By treating contaminant
movement as a ``random walk'' through heterogeneous microstructures, these
models quantify the uncertainty of pollutant attenuation and breakthrough
\cite{Lu2020,Karlen2019}. Integrating AI-driven stochastic simulations enables
researchers to calculate high-fidelity probability distributions for
contaminant spatial distribution, optimizing risk assessment under
high-uncertainty environmental conditions \cite{Li2023,Feng2024}.

Soil aggregates function as high-dimensional structural units where
mineral-organic complexes and microbial consortia regulate fundamental
biogeochemical fluxes \cite{Six2004}. To capture these dynamics, traditional
monitoring is being superseded by AI-driven architectures capable of
processing massive, non-linear datasets
\cite{Lal2015,Baveye2016,Lehmann2020}. Modern assessment utilizes
Convolutional Neural Networks (CNNs) to perform automated feature extraction
on micro-tomographic imagery, quantifying pore-scale connectivity that governs
pollutant transport \cite{Demiral2026,Chen2025}. Furthermore, Random Forest
(RF) and Gradient Boosting Machines (GBM) are deployed to model the
pedotransfer functions of aggregate stability, bypassing the computational
intensity of traditional mechanistic simulations
\cite{DemissieNaveda2026,Guo2026}. By implementing Unsupervised Learning, such
as k-means or DBSCAN, to cluster soil health indicators
\cite{deOliveiraetal2022,Haller2023} and Recurrent Neural Networks (RNNs) to
forecast temporal degradation \cite{Mishra2025,Nunes2025}, researchers can
achieve a high-fidelity, predictive understanding of soil systems that static
environmental frameworks cannot provide.

At the microscale, the interfaces between soil aggregates and transport
pathways exhibit highly nonlinear and dynamic characteristics, reflecting the
complexity of interactions governing pollutant transport, nutrient exchange,
and microbial activity. As a critical component of the biophysical
environment, soil properties are shaped by interconnected geographic and
ecological factors, including geological and geomorphological conditions,
climate, biological activity, and anthropogenic influences
\cite{Brevik2015,Amundson2019}. Together, these factors determine soil
functionality and the provision of ecosystem services
\cite{Baveye2016,Klaucol2013,Klaucol2017,Lemenkova2023}. Consequently, the
ecological significance of soil extends beyond its physicochemical composition
and is intrinsically linked to the interactions among landscape components
and ecosystem processes \cite{Lemenkova2025a}.

%% -------------------------------------------------------------------------
\subsection{Objectives and Scope of the Review}\label{sec1.3}
%% -------------------------------------------------------------------------

This integrative review critically synthesizes more than 100 peer-reviewed
studies to evaluate the efficacy of existing AI and ML methodologies in soil
pollution assessment. While conventional detection techniques often lack the
sensitivity and spatial resolution required for complex matrices, this study
investigates how machine learning architectures---including supervised
classifiers and deep learning models---enhance the characterization of metals,
microplastics, and PFAS.

The primary objective is to transition from descriptive pollutant
identification to predictive computational frameworks. We critically analyse
the performance of RF, GBM, and CNNs in quantifying contaminant binding
mechanisms within soil aggregates. By benchmarking these data-driven
approaches against traditional statistical modelling, this review highlights
the advantages and scalability of AI in ecological monitoring. Ultimately,
this framework integrates remote sensing and automated feature extraction to
provide a unified, high-fidelity methodology for environmental risk assessment
and sustainable soil management.

%% =========================================================================
\section{Material and Methods}\label{sec2}
%% =========================================================================

Speciation and distribution of major pollutants in soil organic complexes
requires a review of their characteristics. Soil contamination is primarily
caused by hazardous substances stemming from both human and natural
activities, Figure~\ref{fig01}.

\begin{figure}[htbp]
    \centering
    \includegraphics[width=\columnwidth]{Fig_01}
    \caption{Soil structure and major sources of its pollution.
        Diagram source: author.}
    \label{fig01}
\end{figure}

%% -------------------------------------------------------------------------
\subsection{Pollutant Groups}\label{sec2.1}
%% -------------------------------------------------------------------------

Major types of soil pollution include the following groups:

\begin{itemize}
    \item \textbf{Heavy metals:} arsenic (As), cadmium (Cd), chromium (Cr),
        mercury (Hg), lead (Pb), cobalt (Co), nickel (Ni), and at higher
        concentrations---copper (Cu), zinc (Zn), iron (Fe).
    \item \textbf{Mineral oils:} petroleum-based hydrocarbons.
    \item \textbf{Organic compounds:} PAHs (Polycyclic Aromatic
        Hydrocarbons), Polychlorinated biphenyl (PCB), and Tributyltin
        (TBT) group.
    \item \textbf{Nutrients:} nitrogen (N), phosphorus (P).
    \item \textbf{POPs} (Persistent Organic Pollutants), including PFAS
        (Per- and Polyfluoroalkyl Substances).
\end{itemize}

Soil contaminant influx originates from synergistic anthropogenic and
lithogenic sources. Anthropogenic forcing---driven by industrial discharge,
mining, intensive agrochemical application, and waste mismanagement---introduces
persistent toxins into the pedosphere \cite{Bwandamuka2025,Latif2026}.
Conversely, natural contributions stem from atmospheric deposition and
seismic-induced geological mobilization \cite{Binetti2026,Latif2026}. These
loading pathways catalyze the hydro-geochemical transfer of contaminants into
groundwater and trophic chains, necessitating advanced modelling to mitigate
systemic ecological and public health risks
\cite{Elmotawakkil2025,Sanad2026}. Recent applications of ML, including
supervised algorithms and CNN, have proven essential for predicting removal
efficiencies and optimizing bioremediation strategies in these complex matrices
\cite{Okafor2026,Pace2026}.

%% -------------------------------------------------------------------------
\subsubsection{Geochemical Origins and Pedospheric Impacts of Heavy Metal
    Contaminants}\label{sec2.1.1}
%% -------------------------------------------------------------------------

The most significant and detrimental type of pollutant affecting soil quality
is heavy metals, as indicated by statistics from the European Environment
Agency (EEA) \cite{EEA2009}. Heavy metals (HM), such as arsenic (As),
cadmium (Cd), chromium (Cr), mercury (Hg), lead (Pb), cobalt (Co), and
nickel (Ni), are toxic trace elements that pose significant risks to soil and
crop growth. These metals are of particular concern due to their harmful
effects on the environment and agriculture. A schematic diagram illustrating
the principal anthropogenic and geogenic pathways of heavy metal influx into
the pedosphere is presented in Figure~\ref{fig02}.

Heavy metal (HM) contamination (e.g., As, Cd, Cr, Hg, Pb) significantly
degrades soil quality and ecosystem services. Originating from industrial
mining \cite{Li2019,Wang2020,AliZerrouki2021,Rasafi2021,EEA2009}, intensive
agriculture, urban emissions \cite{Adewumi2024}, and waste disposal
\cite{LindhLemenkova2022a}, HMs persist and bioaccumulate within the
soil-plant-human axis. Traditional monitoring of these complex biogeochemical
cycles is increasingly replaced by AI and Machine Learning (ML) architectures
to handle high-dimensional, non-linear datasets.

\begin{figure}[htbp]
    \centering
    \includegraphics[width=\columnwidth]{Fig_02}
    \caption{Principal anthropogenic and geogenic pathways of heavy metal
        influx into the pedosphere. Diagram source: author.}
    \label{fig02}
\end{figure}

Modern data-driven analysis utilizes ensemble classifiers and deep learning
to model HM transport and its impact on rhizosphere microbial dynamics. By
integrating geospatial analytics with stochastic simulations, researchers can
automate feature extraction from multisource data, optimizing risk assessment
for leaching and bioaccumulation. These ML-driven pipelines overcome the
limitations of static frameworks, providing high-fidelity, predictive
monitoring of pollutant kinetics across heterogeneous soil matrices.

%% -------------------------------------------------------------------------
\subsubsection{Mobility, Leaching Dynamics, and Ecotoxicological Effects of
    Pesticides}\label{sec2.1.2}
%% -------------------------------------------------------------------------

Pesticide pollution disrupts soil ecosystems through absorption, runoff, and
leaching, impairing biodiversity and nutrient cycling ($N, P, K$),
Figure~\ref{fig03}. Irresponsible application reduces microbial resilience
and facilitates bioaccumulation within the food chain, threatening human
health \cite{GonzalezRodriguez2011,Nauanova2025}. Traditional monitoring of
these dynamics is increasingly augmented by AI and ML frameworks designed for
high-dimensional data-driven analysis. Modern ML architectures, including RF
and CNNs, optimize the prediction of pesticide kinetics and soil organic
matter alterations. By processing multi-temporal datasets, these AI-driven
pipelines quantify non-linear interactions between chemical persistence and
microbial adaptation. Integrating automated feature extraction with geospatial
modelling enables high-fidelity risk assessments of toxicity and leaching.
This shift toward computational pedology allows for real-time predictive
monitoring of pollutant cycles, ensuring the protection of soil fertility and
ecosystem services against chemical stressors.

\begin{figure}[htbp]
    \centering
    \includegraphics[width=\columnwidth]{Fig_03}
    \caption{Principal pathways of pesticide infiltration into the soil
        matrix. Diagram source: author.}
    \label{fig03}
\end{figure}

Leaching facilitates pesticide penetration into groundwater, escalating
ecotoxicological risks and soil structural degradation
\cite{Toumi2025,PerezLucas2025}. This process impairs moisture retention
\cite{Kelsey1968}, diminishes fertility, and disrupts microbial community
dynamics and phosphorus cycling \cite{Ponder2002,Yasir2025}. Resulting
declines in soil biodiversity fundamentally alter ecosystem structures
\cite{Marcelino2024}. To mitigate these risks, AI-driven models are
increasingly deployed to detect and predict leaching patterns. ML
architectures, such as Deep CNN and SVM, utilise data-driven soil analysis
to quantify non-linear interactions between pesticide mobility and
heterogeneous soil matrices. By processing multi-sensor geospatial data,
these AI frameworks provide high-fidelity simulations of contaminant
transport, identifying leaching hotspots that traditional empirical models
overlook. This transition to computational monitoring enables real-time
assessment of groundwater vulnerability, optimizing the preservation of soil
integrity and non-target species.

%% -------------------------------------------------------------------------
\subsubsection{Structural and Biogeochemical Disruption of Soil Systems by
    Microplastics}\label{sec2.1.3}
%% -------------------------------------------------------------------------

Pervasive and persistent pollutants known as microplastics (MP) consist of
materials like polyethylene and polypropylene, and they also serve as
carriers for additional pollutants, particularly heavy metals
\cite{Wang2024a}. A schematic diagram illustrating the major sources of
microplastic contamination of soil is presented in Figure~\ref{fig04}.

\begin{figure}[htbp]
    \centering
    \includegraphics[width=\columnwidth]{Fig_04}
    \caption{Principal pathways of microplastic influx into terrestrial
        soil systems. Diagram source: author.}
    \label{fig04}
\end{figure}

Microplastic (MP) pollution disrupts soil porosity, enzymatic activity, and
hydraulic conductivity, leading to structural degradation and impaired
nutrient cycling, specifically phosphorus uptake
\cite{Akdogan2019,Zhou2024,Wang2022}. These pollutants alter soil aggregate
density and water retention, ultimately stunting plant growth and entering
the food chain \cite{Jing2023,LindhLemenkova2023d,Chenetal2025}. To manage
these complex threats, researchers are pivoting toward AI and ML-driven
analysis. DL models and CNN now automate the detection of MPs within soil
matrices, quantifying non-linear impacts on pore-scale dynamics that
traditional methods miss. These data-driven pipelines integrate multi-source
environmental variables to predict MP transport and bioaccumulation risks,
enabling high-fidelity monitoring of soil health and ecosystem resilience
against plastic-induced stressors \cite{Chen2020,LindhLemenkova2023e}.

%% -------------------------------------------------------------------------
\subsubsection{Environmental Persistence, Speciation, and Transport Pathways
    of PFAS}\label{sec2.1.4}
%% -------------------------------------------------------------------------

Per- and polyfluoroalkyl substances (PFAS) represent a class of anthropogenic
organofluorine compounds engineered to enhance the functional properties of
industrial and consumer products, including fluoropolymer-coated cookware,
food contact materials, and hydrophobic textiles. The exceptional
electronegativity and strength of the $C$--$F$ (carbon-fluorine) bond impart
significant chemical and thermal stability, facilitating the recalcitrance
and long-term persistence of PFAS within environmental matrices, particularly
terrestrial soil profiles. The primary pathways for PFAS infiltration into
soil systems are systematically illustrated in Figure~\ref{fig05}. Modern AI
and ML frameworks are essential for resolving these multifaceted interactions.
Data-driven soil analysis utilises ensemble learning and neural networks to
quantify non-linear PFAS-MP adsorption kinetics and predict leaching risks
\cite{Xu2023}. By integrating deep learning with spectroscopic data, these
ML-driven pipelines automate the identification of PFAS hotspots and simulate
transport dynamics across heterogeneous matrices. This computational shift
enables high-fidelity, real-time monitoring of soil contamination, optimizing
risk assessment under high-uncertainty environmental conditions.

\begin{figure}[htbp]
    \centering
    \includegraphics[width=\columnwidth]{Fig_05}
    \caption{Primary infiltration pathways of PFAS into terrestrial soil
        profiles. Diagram source: author.}
    \label{fig05}
\end{figure}

%% -------------------------------------------------------------------------
\subsubsection{Anthropogenic Nutrient Loading, Soil Degradation, and
    Eutrophication Dynamics}\label{sec2.1.5}
%% -------------------------------------------------------------------------

Anthropogenic hyper-enrichment of primary macronutrients---specifically
nitrogen (N), phosphorus (P), and potassium (K)---is predominantly driven by
the intensive application of synthetic fertilizers and organic manures
\cite{Baghaie2020,Luo2017,Hepperly2009}. These nutrient surpluses induce
significant soil quality degradation and groundwater contamination, ultimately
depressing agricultural productivity. Elevated NPK concentrations further
catalyze deleterious ecological phenomena, including harmful algal blooms and
systemic eutrophication within adjacent aquatic ecosystems
\cite{Pan2019,Guan2024,Xu2025}.

To optimize nutrient management, AI-driven analytical frameworks are
increasingly integrated into soil monitoring. Machine learning (ML)
architectures, including Gradient Boosting and Deep Neural Networks, utilise
data-driven soil analysis to quantify non-linear nutrient leaching kinetics
and forecast spatial and temporal distribution patterns of contaminants. By
leveraging high-resolution remote sensing data and real-time sensor networks,
these ML-driven pipelines automate feature extraction and uncertainty
quantification, allowing for the precise mapping of nutrient flux across
heterogeneous landscape matrices.

PFAS and microplastics (MP) form synergistic contaminants; PFAS adsorption
onto MPs increases pollutant mobility and complicates remediation
\cite{Shen2024}. These substances disrupt carbon (C) and nitrogen (N)
cycling, accelerate litter decomposition, and alter soil pH, ultimately
threatening human health via water and food supplies
\cite{Hossini2025,Appel2022,Cheng2025,Arnesdotter2025}. By processing
multispectral and sensor-derived datasets, these ML pipelines enable
high-fidelity predictive modelling of nutrient cycles, facilitating real-time
interventions to mitigate eutrophication risks. The primary mechanisms of
nutrient infiltration into terrestrial matrices are delineated in
Figure~\ref{fig06}.

\begin{figure}[htbp]
    \centering
    \includegraphics[width=\columnwidth]{Fig_06}
    \caption{Principal mechanisms of nutrient infiltration into the soil
        structure and adjacent aquatic systems. Diagram source: author.}
    \label{fig06}
\end{figure}

Establishing rigorous diagnostic criteria for nutrient enrichment requires a
comprehensive evaluation of concentration gradients relative to soil
texture---specifically the sand, silt, and clay fractions. These granulometric
properties dictate soil porosity and permeability, which govern hydraulic
conductivity and surface runoff potential; for instance, low-permeability
soils accelerate nutrient transport into adjacent water bodies. To resolve
these complex biophysical interactions, AI-driven methods have replaced
traditional linear assessments with high-fidelity, data-driven diagnostics.
Machine Learning (ML) techniques, such as Random Forests (RF) and Support
Vector Machines (SVM), are frequently utilised to model the non-linear
relationship between soil texture and nutrient retention
\cite{Zhang2021,He2022}. Deep Learning (DL) architectures, specifically
Convolutional Neural Networks (CNNs), have been successfully applied to
process hyperspectral imagery for the real-time detection of nitrogen and
phosphorus hotspots across heterogeneous fields \cite{Fanetal2025,Li2023}.
Additionally, Gradient Boosting Machines (GBM) are deployed to predict
leaching risks by integrating multi-source data---including soil moisture,
topography, and fertilization rates---optimizing agricultural management
through predictive uncertainty quantification \cite{Feng2024,Demissie2026}.
Understanding this interplay through AI-enabled pipelines is essential for
mitigating the environmental impact of agricultural nutrient loading.

Nutrients can negatively impact plant development in various ways, primarily
through disrupted growth patterns and stunted root development
\cite{Gang2004,Guo2021}. Overuse of nutrients not only affects soil nutrient
balance but can also lead to alterations in soil pH, which further
complicates plant health \cite{Rahman2014}. The balance of essential minerals
absorbed by plants is crucial, and excess nutrients can hinder their
acquisition, affecting both plant and microbial growth
\cite{Ma2024,Vishwanath2025,Choudhary2022}. This imbalance has downstream
effects, as it impacts microbial activity responsible for the mineralisation
of organic matter and the decomposition of plant litter, contributing to soil
fertility \cite{Atoloye2024}. Additionally, nutrient deposition from
atmospheric sources, such as ammonia~(NH$_{3}$) from cultivated areas, can
harm natural terrestrial habitats, including adjacent forests
\cite{Koukianaki2025}.

%% -------------------------------------------------------------------------
\subsection{Quantitative Data Synthesis and Computational
    Parameterisation}\label{sec2.2}
%% -------------------------------------------------------------------------

Quantitative data synthesis is fundamental to simulating soil response
trajectories under contaminant stress. This process involves the rigorous
parametrization of model dependencies, where empirical results from laboratory
analyses---such as sorption isotherms, hydraulic conductivity, and redox
potential---are translated into high-fidelity mathematical frameworks. The
field has evolved from deterministic equations to AI-driven architectures that
excel in capturing the non-linear, multi-scale interactions within the soil
matrix. Machine Learning (ML) methodologies, specifically Ensemble Learning
and Deep Neural Networks (DNNs), are now utilised to process massive,
heterogeneous datasets that traditional statistical models fail to resolve.
For instance, Random Forest (RF) and Extreme Gradient Boosting (XGBoost)
algorithms are employed to quantify the complex relationship between soil
organic carbon and pollutant sequestration, effectively mapping spatial
heterogeneity in contaminant distribution \cite{Wadoux2024,Sahab2025}.

Furthermore, Convolutional Neural Networks (CNNs) are increasingly applied to
automated feature extraction from micro-CT and hyperspectral imagery, enabling
the characterization of pore-scale connectivity and its impact on solute
transport \cite{Demiraletal2026}. By integrating Transfer Learning and
Physics-Informed Neural Networks (PINNs), researchers can constrain AI models
with known biogeochemical laws, ensuring that predictive monitoring remains
physically consistent even under high-uncertainty environmental conditions
\cite{Karpatne2024,Nguyen2025}. Table~\ref{tab1} delineates the current
state-of-the-art (SOTA) techniques facilitating these advanced soil behaviour
predictions.

\begin{table*}[htbp]
\caption{SOTA techniques and instruments used for detecting pollutant
    concentrations and evaluating the level of soil
    contamination.}\label{tab1}
\begin{tabularx}{\textwidth}{@{} p{0.8cm} p{2.2cm} p{5.5cm} X @{}}
\toprule
\textbf{No} & \textbf{Pollutant} & \textbf{Methods of detection}
    & \textbf{Supporting References} \\
\midrule
1 & HM concentration (Cd, Pb, Me)
    & X-ray fluorescence (XRF) for determining elemental composition of
      soil.
    & \cite{Lu2023,LindhLemenkova2023b,Wan2019,Mukhopadhyay2020,Adewumi2024,Feng2024,Hu2020,Li2019} \\[4pt]
2 & PAH, Tributyltin (TBT)
    & Spectroscopy and X-ray diffraction through Computed Tomography (CT)
      for non-destructive 3-D internal imaging of soil structure.
    & \cite{LindhLemenkova2022d,Deceuster2012,Racic2025,LindhLemenkova2022a,Xu2026,Stojic2019} \\[4pt]
3 & Organic pollutants (PCB)
    & Gas chromatography coupled with mass spectrometry (GC-MS) or
      electron capture detection (GC-ECD).
    & \cite{Bustamante2025,Anagnostou2025,Sandu2025,Shen2024,Appel2022,Arnesdotter2025} \\[4pt]
4 & Soil NPK levels
    & E-sensors, optical transducers, or chemical analysis for evaluation
      of NPK concentration.
    & \cite{Hossain2024,Pal2024,Ma2024,Guan2024,Atoloye2024,Luo2017} \\[4pt]
5 & pH of soil slurry
    & pH-meter or colorimetric test strip; logarithmic scale 0--14
      (neutral = 7; $<$7 acidic; $>$7 alkaline).
    & \cite{Zireeni2023,LindhLemenkova2022b,Rahman2014,Hepperly2009,LindhLemenkova2023c,Guo2021} \\[4pt]
6 & Soil texture
    & Relative \% of sand, silt, and clay by particle size analysis (soil
      texture triangle); e.g., 2.0--0.05\,mm diameter.
    & \cite{Polakowski2023,LindhLemenkova2023d,Mozaffari2026,LindhLemenkova2023a,LindhLemenkova2023e,Jing2023,Wang2022} \\
\bottomrule
\end{tabularx}
\end{table*}

The compiled dataset is synthesized to simulate soil contamination
trajectories, specifically targeting the concentration gradients of heavy
metals, pesticides, microplastics, PFAS, and macronutrients within
ecologically distinct habitats. Quantitative analysis of pollutant--biota
interactions is imperative for characterizing contaminant attributes,
including stoichiometric values and spatial distribution patterns. A primary
bottleneck in environmental data science is the implementation of
reproducible, automated workflows; this is addressed through the deployment
of algorithmic scripts in R or Python
\cite{Olah2025,Lemenkova2022,Wang2024b}, which significantly enhance data
throughput and analytical precision.

Advanced computational modelling is essential for the quantification of
pedological parameters and the visualization of complex biophysical
interactions within soil ecosystems across divergent scenarios, including
anthropogenic climatic shifts, mineral extraction site development, and
strategic urban expansion. By leveraging ML and AI-driven architectures within
Python and R libraries, researchers can execute high-dimensional data analytics
to resolve non-linear correlations between pollutants and soil matrices. These
data-driven pipelines facilitate the generation of robust graphical models and
resistivity visualizations, allowing for a predictive understanding of soil
resilience under varying environmental stressors.

%% =========================================================================
\section{Computational and Spectroscopic Techniques in Soil Pollution
    Detection}\label{sec3}
%% =========================================================================

Current advancements in computational pedology have shifted from traditional
deterministic models toward AI and ML-driven architectures capable of
resolving the high-dimensional, non-linear complexities of soil
aggregate-pollutant interactions. A prominent trend is the deployment of
Convolutional Neural Networks (CNNs) and Deep Learning (DL) for automated
feature extraction from micro-computed tomography ($\mu$-CT) and hyperspectral
imagery, enabling the quantification of pore-scale connectivity and
contaminant localization within mineral-organic complexes
\cite{Demiral2026,Fan2025}.

Ensemble Learning techniques, such as Extreme Gradient Boosting (XGBoost) and
Random Forests (RF), are increasingly utilised to develop robust pedotransfer
functions that predict the sequestration kinetics of heavy metals and PFAS
under high-uncertainty conditions \cite{Sahab2025,Wadoux2024}. To bridge the
gap between empirical data and physical laws, the integration of
Physics-Informed Neural Networks (PINNs) has emerged as a state-of-the-art
approach, constraining stochastic simulations with mass-balance and
thermodynamic principles to ensure biogeochemical consistency in predictive
monitoring \cite{Karpatne2024,Nguyen2025}. These data-driven pipelines, often
implemented via scalable Python or R workflows, facilitate real-time risk
assessment and the identification of leaching hotspots that remain undetected
by standard linear regression frameworks.

\begin{table*}[htbp]
\centering
\caption{Current AI and ML trends in computational methods for soil
    pollutant detection.}
\label{tab:AI_ML_Soil_Trends}
\begin{tabularx}{\textwidth}{p{3.2cm} p{8.5cm} X}
\toprule
\textbf{AI/ML Methodology}
    & \textbf{Application in Soil Aggregates and Pollutants}
    & \textbf{Key Citations} \\
\midrule
Deep Learning (CNNs \& DNNs)
    & Automated feature extraction from micro-CT and hyperspectral imagery
      to quantify pore-scale connectivity and contaminant localization.
    & \cite{Demiral2026,Fan2025} \\[4pt]
Ensemble Learning (XGBoost, RF)
    & Development of robust pedotransfer functions for predicting
      sequestration kinetics and spatial distribution of heavy metals and
      PFAS.
    & \cite{Sahab2025,Wadoux2024} \\[4pt]
Physics-Informed Neural Networks (PINNs)
    & Integrating biogeochemical laws (e.g., mass balance, Darcy's Law)
      into neural networks to ensure physically consistent stochastic
      simulations.
    & \cite{Karpatne2024} \\[4pt]
Transfer Learning
    & Adapting pre-trained models from laboratory-scale experiments to
      predict field-scale pollutant transport and environmental exposure.
    & \cite{Nguyen2025} \\[4pt]
Unsupervised Clustering (DBSCAN, k-means)
    & Identifying multi-temporal spatial heterogeneities and grouping soil
      health indicators to define contamination hotspots.
    & \cite{Haller2023,deOliveiraetal2022} \\
\bottomrule
\end{tabularx}
\end{table*}

%% -------------------------------------------------------------------------
\subsection{Spectroscopic and Synchrotron-Based Characterisation of Soil
    Contaminants}\label{sec3.1}
%% -------------------------------------------------------------------------

Traditional approaches for soil contamination assessment have predominantly
relied on extraction-based analytical methodologies designed to quantify and
characterize the various physicochemical forms of pollutants present within
soil matrices \cite{Sutherland2010}. These conventional methods typically
involve sequential chemical extraction procedures, batch adsorption analyses,
and wet-chemistry techniques aimed at evaluating contaminant mobility,
bioavailability, partitioning behaviour, and potential environmental risks.
Although such approaches remain fundamental in environmental geochemistry and
soil science, their applicability is often constrained by limited spatial
resolution, extensive sample preparation requirements, and difficulties
associated with accurately characterizing complex pollutant interactions within
heterogeneous soil systems.

The methods discussed exhibit several drawbacks primarily linked to the
detection limits of current analytical techniques, which tend to be inadequate
in the context of highly homogeneous samples. Additionally, the adsorption
percentage of contaminants on soil components is often low and uneven, posing
challenges for accurate detection. Proposed solutions to address these issues
include the development of innovative and advanced instruments for assessing
soil quality, alongside the exploration of alternative methods such as remote
sensing \cite{Fan2025,Lemenkova2025b}, machine learning
\cite{Proshad2025,Lemenkova2025d,Lemenkova2024b}, and computer vision
\cite{Fengetal2024}.

In contrast, contemporary analytical frameworks increasingly incorporate
advanced surface-sensitive and high-resolution characterization techniques
capable of investigating contaminant dynamics at micro- and nanoscale levels.
Among these, synchrotron radiation-based methodologies have emerged as highly
effective tools for the detailed investigation of soil contaminants and their
interactions with mineral and organic soil constituents
\cite{Hu2020,Rosskopfova2012,LindhLemenkova2023a}. Synchrotron radiation
technologies employ highly intense and tunable X-ray beams that facilitate
the precise identification of metallic elements, determination of their
oxidation states, and characterization of their chemical speciation within
complex environmental matrices. These techniques are particularly valuable for
elucidating adsorption mechanisms, contaminant transformation pathways, and
interfacial reactions occurring within soil aggregates and porous media.
Furthermore, synchrotron-based analyses significantly improve the capability
to model contaminant distribution patterns and microscale heterogeneity in
contaminated soils.

Additional instrumental techniques representative of advanced analytical
approaches include Fourier transform infrared spectroscopy (FTIR)
\cite{Xu2026} and X-ray photoelectron spectroscopy (XPS)
\cite{Guo2025,LindhLemenkova2023c}. FTIR analysis enables the identification
of molecular functional groups and organic compound interactions within soil
systems, thereby supporting the characterization of pollutant binding
mechanisms and organic matter transformations. Similarly, XPS provides
detailed information regarding elemental composition, surface chemistry, and
electronic states of contaminants and soil particles, allowing for enhanced
interpretation of surface adsorption processes and contaminant speciation.
Collectively, these high-resolution analytical methodologies substantially
improve the detection sensitivity, chemical characterization accuracy, and
mechanistic understanding of pollutants in soil environments.

%% -------------------------------------------------------------------------
\subsection{AI-Augmented Analytical Frameworks for Automated Contaminant
    Detection}\label{sec3.2}
%% -------------------------------------------------------------------------

Recent developments in AI and ML have further transformed soil contamination
studies by enabling the integration and interpretation of large-scale,
multidimensional analytical datasets generated through advanced instrumental
techniques. AI-driven computational frameworks, including artificial neural
networks \cite{Aljanabi2025,Anifowose2024}, random forest algorithms
\cite{Shaheen2018,Foronda2023}, support vector machines
\cite{Agyeman2022,Gholizadeh2020}, and deep learning architectures
\cite{Hu2024,MLReview2024}, are increasingly applied to identify complex
nonlinear relationships among soil physicochemical parameters
\cite{Aljanabi2025,Kammerlander2025}, contaminant distributions
\cite{Agyeman2022,Shaheen2018}, and environmental variables
\cite{Anifowose2024,Gholizadeh2020}. These methodologies facilitate rapid
pattern recognition, predictive modelling, anomaly detection, and spatial
interpolation of contamination data with substantially improved efficiency and
accuracy relative to conventional statistical approaches.

The integration of AI and ML techniques with spectroscopic analyses,
hyperspectral imaging, remote sensing platforms, and geographic information
systems (GIS) has significantly enhanced the capability for real-time
environmental monitoring and high-resolution contamination mapping
\cite{Gholizadeh2020,Hu2024,Lemenkova2025f}. Machine learning-assisted
interpretation of FTIR, XPS, and synchrotron-generated datasets enables
automated feature extraction \cite{Zhou2023,Caporaso2018}, contaminant
classification \cite{Agyeman2022,Hu2024}, and predictive assessment of
pollutant mobility and ecological risk
\cite{Shaheen2018,Anifowose2024,Foronda2023,Gholizadeh2020}. Such
interdisciplinary approaches contribute to the development of precision soil
monitoring systems and advanced decision-support frameworks for sustainable
environmental management and remediation planning.

%% -------------------------------------------------------------------------
\subsection{Computational Frameworks for Pollutant Speciation and
    Geochemical Modelling}\label{sec3.3}
%% -------------------------------------------------------------------------

The methodologies for assessing soil contamination integrated within this
study focus on state-of-the-art (SOTA) computational frameworks designed to
evaluate soil ecosystem services (SES). The systematic workflow for SES
modelling, specifically oriented toward the identification of ecotoxicological
risks and the preservation of soil functional integrity, is delineated in
Table~\ref{tab3}. Traditionally, these methodologies are bifurcated into
empirical models and process-based (mechanistic) models, both of which serve
as fundamental instruments for soil protection, land-use optimization, and
remediation strategy development
\cite{TeranGomez2025,ElSorogy2025,Lemenkova2020,Lemenkova2021,Sandu2025}.

In recent years, the paradigm has shifted toward the integration of ML and AI
to address the inherent non-linearity and spatial heterogeneity of soil
systems. Unlike classical approaches, data-driven ML architectures---including
RF, GBM, and SVR---excel at resolving complex interactions between pedological
variables and anthropogenic pollutants \cite{Li2025a,Racic2025}. For example,
CNNs are increasingly utilised for automated feature extraction from proximal
sensing data to map heavy metal concentrations, while Recurrent Neural
Networks (RNNs) and Long Short-Term Memory (LSTM) networks are deployed to
simulate the temporal dynamics of pesticide leaching and nutrient flux.
Furthermore, the emergence of Physics-Informed Neural Networks (PINNs) allows
for the synthesis of mechanistic hydrological laws with stochastic data
analysis, providing high-fidelity simulations of contaminant transport that
maintain biogeochemical consistency even in data-scarce environments. These
AI-enhanced methodologies significantly improve the precision of SES mapping
and facilitate real-time decision-support systems for environmental monitoring.

\begin{table*}[htbp]
\caption{Current modelling approaches for SES evaluation: characteristics,
    strengths, and limitations.}\label{tab3}
\begin{tabularx}{\textwidth}{@{} p{2.0cm} p{3.5cm} p{3.5cm} p{3.0cm} X @{}}
\toprule
\textbf{Model type} & \textbf{Characteristics} & \textbf{Advantages}
    & \textbf{Limitations} & \textbf{Examples} \\
\midrule
Process-based
    & Use mathematical equations to simulate soil formation and
      physical--chemical processes.
    & Provide high accuracy when well-calibrated with specific biological
      data.
    & Can be overcomplicated due to multiple variables.
    & Integrated Critical Zone (ICZ) model \cite{Giannakis2017} \\[6pt]
Empirical
    & Simple models of soil structure useful for data-scarce areas; rely on
      links between soil and ecosystem services.
    & Efficient for rapid calculation of soil loss (rainfall erodibility,
      slope length, steepness, LULC).
    & Difficulty with dynamic factors (LUC, soil moisture); needs average
      data; prone to overprediction.
    & Universal Soil Loss Equation (USLE) \cite{Benavidez2018} \\[6pt]
Composite
    & Integrate different approaches: process-based $+$ empirical
      frameworks.
    & Provide a holistic view of soil functions and their contribution to
      ecosystem services.
    & Need accurate validation for joint probability distributions.
    & Life Cycle Assessment (LCA) \cite{Sadhukhan2022} \\[6pt]
Statistical
    & Analyse complex links in different SES by capturing non-linear
      relationships.
    & Flexibility for high-dimensional data; accurate detection of tail
      dependencies.
    & Lower accuracy; requires large datasets; structural complexity.
    & Vine copulas \cite{Luetal2020} \\
\bottomrule
\end{tabularx}
\end{table*}

The empirical models yield precise quantification and contribute to the
understanding of soil pollution, although they are challenged by the
ecological soil monitoring limitations associated with the complexity and
variability of different soil types, necessitating further validation. A
prevalent model type for environmental risk assessment estimates ecological
risk factors through contamination indices
\cite{Stojic2019,Mohammad2025,Sun2025,Edo2024,Chon2011,Ferreira2022}. These
methods involve comparing pollutant concentrations in soil samples with
established background levels, where the physical-chemical properties and
contamination levels are examined for ecotoxicity via bioassay techniques. The
summary highlights significant approaches in quantifying ecological indices,
which utilise diverse mathematical frameworks, as outlined in
Table~\ref{tab2_expanded}.

\begin{table*}[htbp]
\centering
\caption{Major ecological risk indices for soil ecosystem assessment.}
\label{tab2_expanded}
\begin{tabularx}{\textwidth}{@{} c p{3.5cm} X p{5.5cm} @{}}
\toprule
\textbf{No} & \textbf{Index} & \textbf{References} & \textbf{Approach} \\
\midrule
1 & ERI -- Ecological Risk Index
    & \cite{Wang2025,Egbueri2020,Adewumi2024,ElSorogy2025,Sidoruk2023,Siddig2025,Fan2025,Proshad2025,AnsoarRodriguez2025}
    & Uses contamination factor ($C_f$) to assess cumulative ecosystem
      risk. \\[4pt]
2 & TRIAD approach
    & \cite{Kim2024,Chon2011,LindhLemenkova2023b,Rasafi2021,Stojic2019,Mohammad2025,Toumi2025,Yasir2025}
    & Integrates chemistry, ecotoxicology, and ecological lines of
      evidence. \\[4pt]
3 & PERI -- Potential Risk Index
    & \cite{AnsoarRodriguez2025,Ferreira2022,Li2025a,Sidoruk2023,Siddig2025,Li2020,Wang2024a,Wang2025}
    & Focuses on metal concentration and biological toxicity
      coefficients. \\[4pt]
4 & $I_{\mathrm{geo}}$ -- Geo-accumulation Index
    & \cite{Li2014,ElSorogy2025,Wang2020,Sidoruk2023,Siddig2025,AliZerrouki2021,LindhLemenkova2022a,Racic2025}
    & Compares current concentrations to lithogenic background
      levels. \\[4pt]
5 & MRI -- Modified PERI by $I_{\mathrm{geo}}$
    & \cite{Liu2021,Ferreira2022,Siddig2025,Li2019,Adewumi2024,Wang2024a,Wang2025,Proshad2025}
    & Enhances accuracy for varied land uses, specifically agricultural
      soils. \\[4pt]
6 & $E_{\mathrm{rf}}$ -- Ecological risk factor
    & \cite{Li2025a,Liu2021,Shetty2025,Marthandan2020,Sun2025,Sidoruk2023}
    & Combines individual concentrations with toxicity and attenuation
      data. \\[4pt]
7 & PLI -- Pollution Load Index
    & \cite{Li2022b,Chukwuonye2025,Wan2019,Mukhopadhyay2020,Sidoruk2023,Li2020,Wang2020,ElSorogy2025}
    & Quantifies total pollution burden relative to baseline
      environments. \\
\bottomrule
\end{tabularx}
\end{table*}

The systematic acquisition of pedological datasets constitutes the critical
foundational phase of soil analysis, requiring rigorous methodological
protocols to ensure the extraction of representative soil aggregate samples.
This multifaceted process initiates with strategic fieldwork designed to
facilitate subsequent high-fidelity laboratory characterization. Essential
parameters of the sampling design include the geospatial demarcation of the
study area and the implementation of optimized sampling geometries---such as
stratified random or systematic grid patterns---utilizing specialized
instrumentation to extract cores across diverse pedogenic horizons and target
depths. To mitigate local stochastic variability, subsamples are homogenized
into composite samples, a procedure vital for stabilizing the signal-to-noise
ratio in downstream analytical workflows.

The challenges and future prospects associated with soil monitoring strategies
are fundamentally oriented toward ensuring the effective functioning, long-term
stability, and protection of soil ecosystems. These integrated monitoring
frameworks incorporate multidisciplinary approaches for assessing soil quality,
contaminant dynamics, ecosystem services, and environmental risks under
changing anthropogenic and climatic pressures. The complexity and
interconnectivity of these processes are schematically illustrated in
Figure~\ref{fig07}, which presents the principal components, analytical
pathways, and monitoring mechanisms involved in comprehensive soil
environmental assessment and sustainable ecosystem management.

\begin{figure}[htbp]
    \centering
    \includegraphics[width=\columnwidth]{Fig_07}
    \caption{Conceptual scheme of SES assessment and monitoring.
        Image source: author.}
    \label{fig07}
\end{figure}

Furthermore, the integration of comprehensive metadata---incorporating
high-resolution meteorological variables and micro-topographic contextual
data---is mandatory for maintaining the geostatistical integrity of the study.
In the contemporary analytical landscape, AI and Machine Learning (ML)
architectures have revolutionised the transition from raw field data to
predictive soil modelling. Traditional sampling efforts are now increasingly
guided by Smart Sampling Strategies powered by Bayesian Optimisation and
Active Learning (AL), which utilise prior environmental covariates to identify
locations of maximum uncertainty, thereby minimising field labour while
maximising information entropy \cite{Wadoux2024}. Moreover, the subsequent
evaluation of biochemical properties, mineralogical assemblages, and multiscale
porosity is significantly enhanced by Deep Learning (DL) pipelines.

For instance, Convolutional Neural Networks (CNNs) are deployed for automated
feature extraction from digital soil morphometrics and X-ray computed
tomography ($\mu$-CT) scans to quantify pore-network tortuosity and aggregate
stability \cite{Demiral2026}. Additionally, Random Forest (RF) and Extreme
Gradient Boosting (XGBoost) algorithms are employed as advanced pedotransfer
functions (PTFs), integrating field metadata with laboratory results to predict
non-linear soil behaviours under varying environmental stressors
\cite{Sahab2025}. These data-driven frameworks facilitate a transition from
static descriptive analysis to dynamic, predictive simulations of soil
ecosystem resilience.

%% =========================================================================
\section{A Multi-Stage Framework for Soil Pollution Monitoring and SES
    Modelling}\label{sec4}
%% =========================================================================

%% -------------------------------------------------------------------------
\subsection{Systematic Workflow for Source Detection and Spatial Hotspot
    Mapping}\label{sec4.1}
%% -------------------------------------------------------------------------

The systematic assessment of soil pollution and its impact on ecosystem
services necessitates an integrated workflow beginning with the identification
of contamination sources through the mapping of hotspot distributions and the
laboratory characterization of physicochemical soil properties, including
heavy metal concentrations, pH, organic matter, and texture,
Table~\ref{tab:ses_workflow_v2}. By applying Principal Component Analysis
(PCA) and related unsupervised machine learning algorithms, researchers can
reduce high-dimensional pollutant datasets to delineate spatial distribution
patterns and elucidate environmental impacts. These insights are subsequently
integrated into empirical and geochemical models---specifically dynamic mass
balance frameworks and AI-enhanced predictive models---to simulate
accumulation and attenuation kinetics within terrestrial systems. Such
modelling enables robust scenario-based forecasting---incorporating optimistic,
default, and pessimistic projections---to predict ecological alterations and
trace the translocation of contaminants from the soil matrix into the trophic
chain.

\begin{table*}[htbp]
\centering
\caption{Workflow summary for soil ecosystem services (SES) modelling and
    risk assessment.}
\label{tab:ses_workflow_v2}
\begin{tabularx}{\textwidth}{p{0.6cm} p{3.0cm} X p{4.5cm}}
\toprule
\textbf{Step} & \textbf{Objective} & \textbf{Methods and Key Components}
    & \textbf{Supporting References} \\
\midrule
1 & Detection of pollution sources
    & Measurement of initial conditions, hotspot mapping via QGIS, lab
      analysis of soil properties (HMs, pH, OM, texture), and PCA for
      spatial distribution assessment.
    & \cite{Adewumi2024,ElSorogy2025,Feng2024,Lemenkova2021,Mukhopadhyay2020,Wan2019} \\[6pt]
2 & Modelling pollutant behaviour
    & Application of empirical, geochemical, and dynamic mass balance models
      to predict concentration trends and simulate accumulation/attenuation
      rates.
    & \cite{Giannakis2017,Koukianaki2025,LindhLemenkova2022b,Mohammad2025,PerezLucas2025,Xu2026} \\[6pt]
3 & Predictive scenario testing
    & Simulation of default, optimistic, and pessimistic scenarios to
      forecast ecosystem changes and modelling pollutant pathways into the
      food chain.
    & \cite{Appel2022,GonzalezRodriguez2011,Hossini2025,Lemenkova2024a,Wang2025,Yasir2025} \\[6pt]
4 & Evaluating ecosystem risk
    & Use of indices ($E_f$, $I_{\mathrm{geo}}$, PERI) to quantify
      pollution levels and assessment of impacts on soil structure,
      porosity, and water capacity.
    & \cite{AnsoarRodriguez2025,Chukwuonye2025,Egbueri2020,Ferreira2022,Kim2024,Li2025a} \\[6pt]
5 & Multi-factor analysis
    & Integration of multiple risk elements for agricultural and forest
      stands; includes Monte Carlo simulations for uncertainty analysis.
    & \cite{Lemenkova2025c,Li2025a,Lu2020,Olah2025,Proshad2025,Wang2024b} \\[6pt]
6 & Prediction and management
    & Remote sensing and time series analysis to forecast long-term impacts
      on health and biodiversity; development of remediation strategies and
      regulatory policies.
    & \cite{Fan2025,Lemenkova2025e,Lemenkova2025d,Pan2019,Stojic2019,Xu2025} \\
\bottomrule
\end{tabularx}
\end{table*}

%% -------------------------------------------------------------------------
\subsection{Ecological Risk Quantification and Scenario-Based Predictive
    Modelling}\label{sec4.2}
%% -------------------------------------------------------------------------

The quantification of soil pollution risks is refined through the application
of ecological indices, including the Enrichment Factor ($E_f$),
Geo-accumulation Index ($I_{\text{geo}}$), and Potential Ecological Risk
Index (PERI). In contemporary research, these indices are often processed via
Artificial Neural Networks (ANNs) or Support Vector Machines (SVMs) to assess
the degradation of soil structural integrity, porosity, and hydraulic capacity
with greater precision than traditional linear methods. To resolve the
multifaceted complexities of contaminant interactions in agricultural and
forest ecosystems, researchers employ multi-factor analysis reinforced by
stochastic Monte Carlo simulations to quantify inherent uncertainties and
parameter sensitivity. Ultimately, the synthesis of multi-sensor remote
sensing data and Long Short-Term Memory (LSTM) time-series analysis
facilitates long-term predictions regarding soil health and biodiversity. This
data-driven computational foundation is essential for formulating targeted
remediation strategies and robust environmental management policies in the
face of anthropogenic stressors.

%% =========================================================================
\section{Discussion}\label{sec5}
%% =========================================================================

Soil contamination constitutes a major ecological and environmental challenge
on a global scale, exerting profound adverse effects on terrestrial
ecosystems, biodiversity, agricultural productivity, and human health. Both
natural environments, including forest ecosystems, and managed agricultural
systems are increasingly threatened by the accumulation of organic and
inorganic pollutants such as heavy metals, pesticides, petroleum hydrocarbons,
industrial waste residues, and emerging contaminants. The effective assessment,
management, and remediation of contaminated soils depend fundamentally on a
comprehensive understanding of the interactions between pollutants and soil
aggregates, which exhibit highly heterogeneous physicochemical properties,
complex microstructures, and variable mineralogical compositions, particle-size
distributions, and organic matter contents. These intrinsic soil
characteristics strongly influence contaminant mobility, bioavailability,
persistence, and toxicity within the soil matrix.

This study critically reviews current methodologies employed for soil pollution
detection and characterization, including conventional physicochemical
analyses, spectroscopic techniques, geostatistical approaches, and emerging
AI-assisted diagnostic systems. In addition, it proposes strategic and
multidisciplinary frameworks for the prevention, monitoring, and remediation
of soil contamination, emphasizing the importance of sustainable management
practices, advanced computational modelling, and data-driven environmental
governance to mitigate long-term ecological and agricultural impacts.

Recent advances in analytical technologies, remote sensing, and geospatial
monitoring have significantly improved the detection and characterization of
soil pollution. In particular, the integration of artificial intelligence (AI)
and machine learning (ML) approaches has emerged as a transformative
development in soil contamination assessment. AI- and ML-based models enable
the rapid processing of large and multidimensional environmental datasets,
facilitating the prediction, classification, and spatial mapping of
contaminated sites with enhanced precision and efficiency. Techniques such as
artificial neural networks, support vector machines, random forests, deep
learning architectures, and hybrid predictive frameworks have demonstrated
substantial potential for identifying contamination patterns, estimating
pollutant concentrations, and optimizing remediation strategies. Furthermore,
the incorporation of AI-driven data analytics with remote sensing platforms,
hyperspectral imaging, and Internet of Things (IoT)-based soil sensors has
improved real-time monitoring capabilities and strengthened decision-making
processes in environmental risk assessment.

The analysis of the scientific literature indicates that contemporary research
on soil contamination is predominantly focused on three principal objectives.
The first involves the identification, characterization, and tracing of
pollution sources in contaminated soils, including anthropogenic inputs derived
from industrial activities, agricultural intensification, mining operations,
urbanization, and waste disposal practices. The second objective concerns the
evaluation of pollutant adsorption mechanisms, retention capacities, mobility
patterns, and spatial distribution within heterogeneous soil matrices.
Particular attention has been devoted to understanding the physicochemical
interactions between contaminants and soil constituents, including clay
minerals, organic matter fractions, oxides, and microbial communities, which
collectively govern contaminant transport, persistence, and bioavailability.
The third major research direction examines the phyto-uptake of pollutants by
plants, the transfer of contaminants through trophic pathways, and the
associated ecotoxicological impacts on terrestrial ecosystems, crop
productivity, food safety, and human health.

Despite substantial scientific progress, the literature demonstrates that only
a limited number of studies conducted over the past decade have provided
comprehensive scoping reviews addressing soil pollution within integrated
soil-plant systems. Existing investigations often remain fragmented, focusing
on isolated contaminants, specific remediation technologies, or localized
environmental conditions, thereby limiting the development of holistic and
standardized assessment frameworks. Consequently, there remains a critical need
for multidisciplinary and large-scale review studies capable of synthesizing
advances in soil chemistry, ecotoxicology, environmental engineering,
computational modelling, and AI-driven analytical techniques.

%% =========================================================================
\section{Conclusions}\label{sec6}
%% =========================================================================

In this study, a comprehensive evaluation of multiple sources of soil
contamination was conducted, with particular emphasis on toxic metals,
pesticides, microplastics, per- and polyfluoroalkyl substances (PFAS), and
excessive nutrient accumulation. These contaminants represent major
environmental stressors due to their persistence, mobility, bioaccumulation
potential, and adverse impacts on soil quality, ecosystem stability,
agricultural productivity, and human health. The investigation systematically
examines the physicochemical behaviour, environmental pathways, and ecological
implications of these pollutant groups within diverse soil systems.

In recent years, the incorporation of artificial intelligence (AI) and machine
learning (ML) methodologies into soil-plant system studies has considerably
expanded the analytical capabilities available for contamination assessment
and environmental prediction. Advanced ML algorithms, including random forest
models, convolutional neural networks, support vector machines, and ensemble
learning techniques, are increasingly utilized to predict contaminant
dispersion, assess ecological risks, model plant uptake behaviour, and identify
complex nonlinear relationships among soil physicochemical parameters and
pollutant dynamics. Moreover, AI-assisted integration of remote sensing data,
hyperspectral imaging, geostatistical analysis, and sensor-based monitoring
systems has improved the spatial and temporal resolution of contamination
mapping, thereby supporting more accurate environmental diagnostics and
remediation planning.

The study further provides a detailed assessment of current identification
techniques and analytical approaches employed for contaminant detection,
characterization, and monitoring. Particular attention is devoted to
conventional laboratory-based methodologies, including chromatographic,
spectroscopic, and geochemical analyses, alongside emerging technologies such
as hyperspectral imaging, remote sensing, geospatial analysis, and
sensor-based monitoring systems. In addition, the paper highlights the
methodological distinctions among existing computational and environmental
modelling tools used in soil ecosystem analysis, emphasizing their respective
capabilities, limitations, spatial resolutions, predictive accuracies, and
applicability under varying environmental conditions.

A central component of this research involves the integration of artificial
intelligence (AI) and machine learning (ML) techniques into soil contamination
assessment frameworks. Advanced computational approaches, including artificial
neural networks, random forest algorithms, support vector machines, gradient
boosting models, and deep learning architectures, are critically evaluated for
their effectiveness in predicting contaminant distribution, estimating
concentration levels, identifying pollution hotspots, and tracing pollutant
origins. These AI-driven methodologies facilitate the processing of large-scale
and multidimensional environmental datasets, enabling improved pattern
recognition, uncertainty reduction, and predictive performance compared with
conventional statistical methods. Furthermore, the combination of ML algorithms
with remote sensing technologies and geographic information systems (GIS)
significantly enhances the accuracy of spatial contamination mapping and
environmental risk assessment.

The study proposes strategic solutions for the development of robust and
adaptive modelling frameworks capable of accurately identifying and quantifying
soil contaminant concentrations while simultaneously determining the primary
anthropogenic and natural sources of pollution. These proposed frameworks
emphasize interdisciplinary integration among soil science, environmental
chemistry, ecotoxicology, data science, and computational modelling in order
to improve the reliability and scalability of contamination assessment
methodologies. Particular consideration is given to the implementation of
explainable AI approaches, hybrid modelling systems, and data fusion strategies
to support transparent and scientifically interpretable decision-making
processes.

Additionally, the project interface incorporates a broad interdisciplinary
analysis of soil environments supported by a comprehensive review of more than
100 scientific sources indexed in major bibliographic databases, including
Scopus and Web of Science. This extensive literature synthesis provides a
robust scientific foundation for evaluating recent advances in soil
contamination research and emerging technological applications. Overall, the
study contributes substantially to the advancement of methodological frameworks
for the detection, mapping, characterization, and predictive assessment of soil
contaminants, while reinforcing the role of AI- and ML-assisted systems in the
future development of precision environmental monitoring and sustainable soil
management practices.

Major modelling trends associated with the evaluation of soil contamination by
per- and polyfluoroalkyl substances (PFAS) were comprehensively analysed, with
particular emphasis on their methodological capabilities, predictive
performance, and inherent limitations across diverse environmental
applications. PFAS contamination presents a significant challenge due to the
exceptional chemical stability, environmental persistence, mobility, and
bioaccumulative properties of these compounds, which complicate their
detection, transport modelling, and remediation within terrestrial ecosystems.
The review critically examined a broad spectrum of modelling approaches,
including physicochemical transport models, geostatistical frameworks,
fate-and-transport simulations, remote sensing applications, and data-driven
computational methodologies designed to assess PFAS distribution, migration
pathways, and ecological risks in soil systems.

The analysis demonstrated that contemporary soil monitoring methodologies
substantially reinforce the conventional framework of soil ecosystem services
(SES) by enabling more precise characterization of soil functionality,
environmental quality, and ecosystem resilience. This advancement is primarily
achieved through the integration of multi-source and high-dimensional datasets
derived from laboratory analyses, hyperspectral imaging, satellite remote
sensing, geospatial technologies, environmental sensors, and long-term
monitoring networks. Such integrated datasets facilitate robust environmental
evaluations through the identification of spectral fingerprints associated with
soil components, contaminant signatures, mineralogical properties, and organic
matter dynamics. These spectral and geospatial approaches significantly improve
the detection sensitivity and spatial resolution of contamination assessments,
particularly in heterogeneous and large-scale soil environments.

Furthermore, the review highlights the growing role of artificial intelligence
(AI) and machine learning (ML) in advancing PFAS contamination modelling and
soil ecosystem analysis. AI-driven predictive systems, including deep learning
models, convolutional neural networks, ensemble learning techniques, and hybrid
geospatial algorithms, have increasingly demonstrated strong potential for
identifying complex nonlinear relationships between PFAS behaviour and soil
physicochemical parameters. Machine learning approaches also contribute to
improved prediction of contaminant transport, adsorption behaviour, leaching
potential, and long-term environmental persistence under variable climatic and
land-use conditions. In addition, AI-assisted integration of hyperspectral
data, remote sensing imagery, and geographic information systems (GIS) enables
high-resolution spatial mapping and real-time environmental diagnostics,
thereby enhancing both monitoring efficiency and risk assessment accuracy.

The review additionally examined innovative agronomic and environmental
management techniques aimed at improving soil ecosystem services within
contaminated and vulnerable agricultural systems. These approaches include
precision agriculture technologies, sustainable nutrient management,
phytoremediation strategies, biochar applications, regenerative farming
practices, and sensor-based soil health monitoring systems. Particular
attention was devoted to the evaluation of performance indicators and
quantitative metrics associated with soil ecosystem services, including
nutrient cycling capacity, carbon sequestration potential, water retention
efficiency, biodiversity preservation, and contaminant buffering functions.
The integration of AI-supported analytical tools into these management
frameworks further enhances the capacity to optimize agricultural practices
while minimizing environmental degradation and maintaining long-term soil
productivity.

Overall, the review underscores the importance of interdisciplinary and
data-driven methodologies for improving the assessment and management of
PFAS-contaminated soils. It also highlights the critical need for standardized
modelling frameworks, enhanced monitoring infrastructures, and advanced
AI-enabled analytical systems capable of supporting sustainable soil
management, environmental protection, and evidence-based policy development in
response to increasing global contamination pressures.

%% =========================================================================
%% Declaration of competing interest
%% =========================================================================
\section*{Declaration of Competing Interest}
The author declares that there is no conflict of interest.

%% =========================================================================
%% CRediT authorship contribution statement
%% =========================================================================
\section*{CRediT Authorship Contribution Statement}
\textbf{Polina Lemenkova:} Conceptualization, Methodology, Writing ---
original draft, Writing --- review \& editing, Visualization.

%% =========================================================================
%% Acknowledgements
%% =========================================================================
\section*{Acknowledgements}
The author acknowledges support from the Bureau de Recherches G\'eologiques
et Mini\`eres (BRGM) and LE STUDIUM Institut d'Etudes Avanc\'ees du Val de
Loire.

%% =========================================================================
%% Bibliography
%% =========================================================================
\bibliographystyle{elsarticle-num}
\bibliography{Bibliography}

\end{document}
